\section{Discretisation numerique}

% Du modele continu au schema numerique.

\subsection{Schema d'Euler semi-implicite}

% Pourquoi semi-implicite plutot qu'explicite : meilleur comportement en temps
% long sur les systemes mecaniques, et c'est lui qui rend possible la
% formulation par minimisation.

\subsection{Formulation par minimisation}

% Un pas de temps vu comme la minimisation d'un compromis entre inertie et
% energie potentielle.

\subsection{Prise en compte des contraintes : une projection}

% Minimiser sous contrainte revient a projeter la position libre sur
% l'ensemble admissible.

\subsection{Linearisation et projection iterative}

% L'ensemble admissible n'etant pas convexe, on linearise autour de la
% configuration courante. Decrire l'algorithme et son critere d'arret.

\subsection{Relecture de la vitesse}

% $\vit^{n+1} = (\pos^{n+1} - \pos^n)/\dt$. C'est la que le choc se produit :
% la vitesse n'est pas transportee, elle est deduite du deplacement reellement
% effectue. Insister, c'est le coeur du schema.

\subsection{Choix du pas de temps}

% Critere : le deplacement par pas doit rester petit devant le rayon.
